\documentclass[11pt]{article}

\usepackage[margin=1in]{geometry}
\usepackage{setspace}
\usepackage{hyperref}

\title{Renovating Calculus: A Summary of a Case Study from Augsburg University}
\date{}

\begin{document}

\maketitle

\onehalfspacing

\section*{Introduction}

Across many universities, faculty are rethinking how calculus is taught in order to better support student success and engagement. This document summarizes a case study from Augsburg University \href{https://arxiv.org/pdf/2406.08508}{Download the paper} that describes how a mathematics department redesigned its calculus sequence through interdisciplinary collaboration and a focus on relevance, conceptual understanding, and transfer of knowledge. The goal of this summary is to highlight the key motivations, changes, and outcomes of that redesign and to provide a starting point for discussion among instructors.

\section*{Institutional Context}

Augsburg University is a small private institution in Minneapolis, Minnesota with a long history of serving a diverse student population. The university enrolls many first-generation students, adult learners, and students from immigrant communities. It also supports students with disabilities and students recovering from chemical dependency through specialized programs. Over the past decade the demographic composition of Augsburg’s student body has changed significantly, with a substantial increase in the number of students identifying as people of color.

These changes prompted faculty to reflect on how the curriculum could better serve the needs and interests of their students. At the institutional level, Augsburg supports ongoing curricular revision, and work on teaching and learning is recognized as Scholarship of Teaching and Learning within expectations for promotion and tenure. Within the Mathematics, Statistics, and Computer Science department, faculty collaboration is common, and instructors regularly work with colleagues in disciplines such as biology, chemistry, economics, and business.

Because Augsburg is a relatively small institution, faculty have considerable flexibility to experiment with new pedagogical approaches. Courses often incorporate active learning and inquiry-based learning, and the department has a strong culture of collaboration and peer support. Faculty members also recognize that many students benefit from additional academic support, which has encouraged the development of courses designed to engage students and help them succeed.

\section*{The Structure of Calculus at Augsburg}

Augsburg offers a three-course calculus sequence consisting of Calculus I, Calculus II, and Multivariable Calculus. Students taking these courses typically major in fields such as biology, chemistry, data science, computer science, economics, education, mathematics, or physics.

Each calculus course meets three times per week for lecture and includes an additional weekly laboratory session. Prior to the curricular redesign, most lecture sessions were delivered using traditional lectures supplemented with occasional group activities. The laboratory sessions often focused on applications or on computational tools such as Excel or Mathematica.

Over several decades, the department had gradually shifted the curriculum toward more applied content. For example, some integration techniques had already been replaced with introductory work on differential equations. However, faculty recognized that further change was necessary to keep pace with the evolving needs of their students.

\section*{Goals of the Curricular Renovation}

Augsburg participated in the SUMMIT-P project, a national initiative supported by the National Science Foundation that focuses on improving undergraduate mathematics education through interdisciplinary collaboration.

The Augsburg team framed their work not as a complete overhaul of the calculus curriculum but as a renovation. This framing was intentional. It acknowledged that the existing courses were already successful in many ways while also creating space for improvement. Pass rates in the calculus sequence were already relatively high, and the goal was to enhance the curriculum rather than replace it.

Through discussions within the department and with faculty from other disciplines, the team developed a set of shared goals for the calculus sequence. They hoped students would:

\begin{itemize}
\item Recognize that calculus is useful and relevant.
\item Be able to apply calculus concepts in new contexts.
\item Develop both conceptual understanding and procedural fluency.
\item Feel encouraged to continue studying mathematics.
\end{itemize}

These goals emphasized not only mastery of technical skills but also students’ ability to use mathematics meaningfully in other fields.

\section*{Streamlining the Curriculum}

One of the first challenges the team encountered was the sheer amount of content already present in the calculus sequence. Faculty wanted to create more opportunities for applications and interdisciplinary connections, but doing so required reducing some existing topics.

To guide these decisions, the team organized a series of discussions with faculty from partner disciplines, including biology, environmental science, chemistry, economics, business, physics, and computer science. These conversations revealed which calculus topics were most important for students in those fields.

Interestingly, the discussions showed that many partner disciplines rely primarily on relatively simple mathematical ideas rather than the more elaborate techniques often emphasized in calculus courses. For example, disciplines frequently use functions such as polynomials and exponentials but rarely require complicated combinations of functions.

Similarly, several integration techniques traditionally taught in calculus were not used in subsequent courses in other departments. As a result, topics such as partial fractions and trigonometric substitution were removed from the curriculum.

By eliminating or reducing some topics and moving others between courses, faculty created space to introduce ideas that were more relevant to students’ future studies. For instance, differential equations now appear earlier in the sequence, and Multivariable Calculus has been made more accessible to students who complete only two semesters of calculus.

\section*{Core Ideas in Calculus I}

Because many students take only a single semester of calculus, the faculty decided to center Calculus I around several fundamental ideas that they felt every student should encounter. These ideas include:

\begin{itemize}
\item Rate of change and differentiation
\item The concept of a limit
\item Calculus as a tool for approximation
\item Optimization
\item The Fundamental Theorem of Calculus
\end{itemize}

This approach ensures that even students who do not continue beyond the first semester encounter the central concepts that make calculus useful in many disciplines.

\section*{Making Calculus Relevant}

A major goal of the redesign was to help students see the relevance of calculus. Research in mathematics education suggests that students are more engaged and motivated when they see how mathematical ideas connect to real-world situations.

To support this goal, the Augsburg faculty incorporated applied examples into nearly every class meeting. Instead of presenting theory first and applications later, instructors often begin with a real-world situation and use it to motivate the mathematics that follows.

Students regularly work in small groups on exploratory activities during class. These activities draw on examples from many fields, including biology, economics, physics, engineering, medicine, and environmental science. Faculty also designed examples that connect to everyday experiences, such as driving, weather, or alcohol metabolism.

Interdisciplinary collaboration played an important role in developing these activities. Faculty from other departments helped identify authentic examples and ensured that the mathematical models used in class reflected how mathematics is actually applied in those disciplines.

\section*{Encouraging Transfer of Learning}

Another major focus of the redesign was helping students transfer their mathematical knowledge to new contexts. Many students learn procedures in mathematics courses but struggle to apply those ideas elsewhere.

To address this challenge, Augsburg uses weekly laboratory sessions to emphasize transfer of learning. During these sessions, students work in small groups on extended applied problems drawn from various disciplines.

A distinctive feature of these activities is that they intentionally require students to use mathematical ideas from earlier in the course rather than the most recently learned topic. This encourages students to recognize when previously learned concepts are relevant and to practice retrieving and applying them independently.

The problems are structured so that students must analyze the situation, interpret data, and determine which mathematical tools are appropriate. Instructors guide the process by asking questions and encouraging reflection, but they avoid explicitly telling students which techniques to use.

\section*{Outcomes}

After several years of implementing the redesigned curriculum, faculty observed several encouraging outcomes. Students reported that they found the applications interesting and relevant to their studies. Many also expressed a stronger sense of belonging and confidence in their ability to succeed in mathematics.

Instructors noticed that students were more engaged during class and more willing to collaborate with their peers. Perhaps most strikingly, students rarely asked the familiar question, “When will we ever use this?” Instead, they increasingly recognized the usefulness of calculus for understanding problems in a wide range of disciplines.

Overall, the Augsburg experience suggests that focusing on core ideas, emphasizing applications, and encouraging interdisciplinary collaboration can help make calculus more meaningful and accessible for students.

\end{document}